\subsection{Baselines and Experimental Setup}
\label{sec:baselines}

To demonstrate the superior performance and robustness of our proposed CLIP-Enhanced Semantic Manifold Diffusion (CESM-Diff) model, we perform a comprehensive comparison against a diverse collection of state-of-the-art (SOTA) methods. These baselines are selected from three distinct yet related paradigms: self-supervised contrastive learning frameworks, advanced generative models, and specific zero-shot fault diagnosis techniques. This multi-faceted comparison allows us to evaluate the efficiency of our manifold-constrained diffusion approach from multiple perspectives, including representation quality, generative fidelity, and zero-shot generalization capability.

\subsubsection{Self-Supervised Contrastive Learning Paradigm}
We include several prominent self-supervised paradigms to assess the inherent benefit of our semantic manifold approach against methods that learn representations purely from data augmentations. Deep learning~\cite{lecun2015deep} has fundamentally transformed feature extraction for industrial signal processing; however, the degree to which self-supervised objectives alone can replace semantically grounded generation remains an open question.
\begin{itemize}
    \item \textbf{MoCo-v3} \cite{chen2021mocov3}: Employed as a powerful contrastive learning baseline that utilizes a momentum encoder. MoCo-v3 is highly effective at learning view-invariant features; comparing it with CESM-Diff helps us evaluate if simple instance-level contrastive learning can match the performance of explicit semantic-to-feature mapping.
    \item \textbf{DINO} \cite{caron2021dino}: Represents a self-distillation approach that operates without explicit labels, focusing on local-to-global matching to learn patch-level and signal-level features simultaneously. We adapt this for 1D vibration signals using window-based local cropping and global alignment.
    \item \textbf{Barlow Twins} \cite{zbontar2021barlow}: Avoids representation collapse by ensuring the cross-correlation matrix between distorted versions of the same sample is close to the identity matrix. This is used to test the stability of feature decorrelation and non-redundancy in high-dimensional industrial sensor data.
    \item \textbf{VICReg} \cite{bardes2021vicreg}: Employs Variance-Invariance-Covariance regularization to learn informative embeddings without the necessity of using negative pairs. It provides a strong benchmark for non-contrastive self-supervision in industrial signal processing.
\end{itemize}

\subsubsection{Generative and Zero-Shot Fault Diagnosis Models}
For a direct comparison in the specific context of zero-shot fault diagnosis (ZSFD), where labels for test classes are entirely withheld, we compare our method against specialized generative architectures designed to synthesize unseen class features. Our diffusion-based generation builds upon the foundational denoising diffusion probabilistic model (DDPM)~\cite{ho2020denoising}, which has demonstrated remarkable sample quality through iterative Gaussian denoising steps. Additionally, we leverage the vision-language alignment capabilities of CLIP~\cite{radford2021learning} to provide semantically rich conditioning signals that guide the diffusion process on the learned manifold.
\begin{itemize}
    \item \textbf{DiffusionSSL} \cite{li2023diffusionssl}: A state-of-the-art diffusion-based self-supervised learner that generates high-fidelity synthetic samples to proactively bridge the semantic gap between seen and unseen fault classes. This serves as the primary diffusion-based sibling to our framework.
    \item \textbf{DiffuCDD}: Serves as a specialized generative baseline for cross-domain diagnosis using conditional diffusion foundations. It represents current industry best practices for industrial sensor data augmentation under limited label availability.
    \item \textbf{f-CLSWGAN}: A classic zero-shot learning method that combines Wasserstein GANs with a classification consistency loss to generate discriminative features for unseen classes based on static attribute vectors. It provides a baseline for traditional GAN-based ZSL approaches.
    \item \textbf{CVAE (Conditional VAE)}: A probabilistic generative model grounded in the variational autoencoder framework~\cite{kingma2013auto}, used to compare our manifold-constrained diffusion against traditional density-estimation-based feature generation. CVAE benchmarks the performance of learned latent distributions under Gaussian assumptions.
\end{itemize}

For all generative baselines, the underlying feature extractor employs a ResNet-18 backbone~\cite{he2016deep} pre-trained on ImageNet and subsequently fine-tuned on the seen-class training data. This ensures that every compared method benefits from the same high-quality deep residual representations as a starting point, isolating the contribution of the generative mechanism itself.

Unlike these existing methods that often rely on static or unconstrained latent spaces, our \textbf{CESM-Diff (Ours)} leverages the CLIP-driven semantic manifold to perform the diffusion process directly in a geometrically informed and linguistically grounded space. By aligning physical sensor data with high-level textual descriptions through the manifold-constrained diffusion process, our model aims to capture the fine-grained nuances of unseen fault modes more effectively than standard contrastive or unconstrained generative approaches.

\subsubsection{Evaluation Metrics}
We adopt two complementary metrics to comprehensively evaluate zero-shot fault diagnosis performance:
\begin{itemize}
    \item \textbf{Mean Class Accuracy (MCA)}: Defined as $\mathrm{MCA} = \frac{1}{C} \sum_{c=1}^{C} \frac{n_c^{\text{correct}}}{n_c}$, where $C$ is the total number of classes, $n_c$ is the number of test samples in class $c$, and $n_c^{\text{correct}}$ is the number of correctly classified samples in class $c$. MCA prevents majority classes from dominating the evaluation, which is critical in fault diagnosis where class distributions are often highly imbalanced.
    \item \textbf{Harmonic Mean Accuracy (H-mean)}: The harmonic mean of seen-class accuracy $A_s$ and unseen-class accuracy $A_u$, computed as $\mathrm{H\text{-}mean} = \frac{2 \cdot A_s \cdot A_u}{A_s + A_u}$. H-mean penalizes models that achieve high accuracy on one group at the expense of the other, making it particularly useful for generalized zero-shot learning (GZSL) settings where both seen and unseen classes appear at test time.
\end{itemize}

\subsubsection{Cross-Validation and Statistical Robustness}
To ensure statistical reliability, we employ a stratified $k$-fold cross-validation protocol within the seen-class partition. Specifically, we perform 5-fold stratified cross-validation during hyperparameter selection: the seen-class training data is split into five equal folds while preserving the original class distribution, with four folds used for training and the remaining fold reserved for validation. Hyperparameters---including the learning rate, the balance weights $\lambda_{acl}$ and $\lambda_{smi}$, and the diffusion noise schedule---are tuned exclusively on these validation folds to prevent any information leakage from unseen classes. After the optimal configuration is determined, the final model is retrained on the complete seen-class training set and evaluated on the held-out unseen classes. Each experiment is repeated 10 times with different random seeds, and we report the mean and standard deviation to account for stochastic variation in both the diffusion sampling process and the random initialization of learnable parameters.

\subsubsection{Detailed Training Configurations and Fairness}
All baselines are implemented using their optimal hyper-parameter configurations as reported in their respective original literature, carefully adapted to the three industrial datasets described in Section~\ref{sec:datasets}. For the self-supervised methods (MoCo, DINO, etc.), we append a linear probing layer or a nearest-centroid classifier to evaluate their zero-shot generalization capabilities based on the CLIP-derived class prototypes. We ensure that the training conditions---including batch size (256), training epochs (500), and hardware resources---are strictly identical for all methods to guarantee a fair and unbiased comparison of their diagnostic efficacy across the TEP, Hydraulic System, and CWRU datasets. This controlled environment ensures that any observed performance gains are due to the architectural innovations of CESM-Diff rather than variations in training intensity or hardware acceleration.
